% Generated draft: requires mathematical and source review.
\documentclass[11pt]{article}
\usepackage[margin=25mm]{geometry}
\usepackage{fontspec}
\setmainfont{Latin Modern Roman}
\usepackage{amsmath,amssymb,mathtools,unicode-math}
\setmathfont{Latin Modern Math}
\usepackage{enumitem,needspace,xurl,hyperref}
\setlength{\parindent}{0pt}
\setlength{\parskip}{5pt}
\setlength{\emergencystretch}{3em}
\newcommand{\N}{\mathbb{N}}
\newcommand{\Z}{\mathbb{Z}}
\newcommand{\Q}{\mathbb{Q}}
\newcommand{\R}{\mathbb{R}}
\newcommand{\C}{\mathbb{C}}
\newcommand{\F}{\mathbb{F}}
\newcommand{\A}{\mathbb{A}}
\newcommand{\PP}{\mathbb{P}}
\newcommand{\E}{\mathbb{E}}
\newcommand{\eps}{\varepsilon}
\newcommand{\dd}{\,\mathrm{d}}
\begin{document}
\section*{UnsolvedMath ID: 30003105}
\textbf{Source problem code:} catalogue entry
\subsection{Definitions and mathematical statement}
Let \(p\) be a prime and let \(n\geq 1\). Write \(W_{\Q_p}\) for the Weil group of \(\Q_p\), and let \(\operatorname{Bun}_{GL_n}\) denote the stack of rank-\(n\) vector bundles in the intended Fargues--Fontaine setting. The source does not specify the coefficient field of the parameter, the precise geometric base, or the category of sheaves.

\textbf{Source question, with unresolved conventions.} For every discrete Weil parameter
\[
\varphi:W_{\Q_p}\longrightarrow GL_n,
\]
does there exist a sheaf \(F_{\varphi}\) on \(\operatorname{Bun}_{GL_n}\), up to shifts and twists, satisfying both of the following requirements?
\begin{enumerate}
\item Its Hecke transform has eigenvalue \(\varphi\), expressed schematically in the supplied record as
\[
\operatorname{Hecke}(F_{\varphi})\simeq F_{\varphi}\otimes\varphi.
\]
\item Its pullbacks to the basic points correspond to the local Langlands representations attached to \(\varphi\).
\end{enumerate}
Here the displayed Hecke identity preserves the source formulation: it is not yet a typed isomorphism in a specified sheaf category. The source does not identify the Hecke functor, its target, the geometric realization of the eigenvalue, or the required compatibility isomorphisms. It also does not specify the permitted shifts and twists.

The term ``discrete Weil parameter'' requires clarification. The record does not specify continuity or semisimplicity conventions, the meaning of discreteness, or whether a Weil--Deligne parameter with monodromy, or an equivalent enlarged parameter, was intended. These distinctions affect the scope of the question and cannot be resolved by silently replacing the displayed parameter.

Likewise, the basic-point requirement needs the relevant basic loci or residual gerbes, their automorphism groups, the associated local Langlands representations, and the normalization relating representations to sheaf restrictions. A pullback to a geometric point alone does not specify the asserted representation-theoretic correspondence. Thus the record supplies an existence question, but these missing conventions prevent a fully precise formulation.
\subsection{Short English statement}
The question asks whether every discrete parameter in the supplied sense admits a Fargues--Fontaine Hecke eigensheaf for \(GL_n\), up to shifts and twists, whose restrictions at basic points recover the associated local Langlands representations. The parameter convention, sheaf category, Hecke operation, and restriction normalizations remain unspecified. The supplied partially solved label and literature assessment are dated curation, not independently verified conclusions.
\subsection{Sources}
\begin{itemize}
\item Supplied local UnsolvedMath registry record: permanent ID \(30003105\), source problem code \texttt{catalogue entry}, title ``Fargues--Fontaine Hecke Eigensheaves for General Linear Groups.'' The record attributes the question to Oberwolfach Reports workshop discussions; no individual proposer is supplied. Its proposed year is 2016.
\item Supplied source citation: \emph{Arbeitsgemeinschaft: The Geometric Langlands Conjecture} (2017), Oberwolfach Reports, pp.~1027--1098, DOI \texttt{10.4171/owr/2016/20}. \url{https://doi.org/10.4171/owr/2016/20}. The supplied citation year is 2017, while the DOI contains 2016; their relationship has not been independently checked.
\item Supplied literature source, ``OWR report'': \url{https://ems.press/content/serial-article-files/46624}.
\item Supplied literature source, ``Anschuetz--Le Bras'': \url{https://www.cambridge.org/core/journals/journal-of-the-institute-of-mathematics-of-jussieu/article/averaging-functors-in-fargues-program-for-mathrm-gln/471BA6B4BB71261287EC350F1C8DECA2}.
\item Supplied literature source, ``Fargues--Scholze'': \url{https://arxiv.org/abs/2102.13459}.
\item Literature checked at: \texttt{2026-08-21}. This is supplied dated curation, not an independently verified assessment. The record labels the problem partially solved and states that the \(GL_n\) cuspidal or irreducible case is proved by Anschuetz--Le Bras and Hansen, while the full discrete-parameter statement represented here is not established. The record explicitly requests independent verification of time-sensitive status and cited claims. No separate Hansen reference is supplied.
\item The supplied statement status is \texttt{corrected\_verified}; its accompanying explanation says only that notation, typography, and sentence structure were repaired without changing mathematical content. This does not establish mathematical correctness or resolve the formulation issues above.
\end{itemize}
\end{document}
